\documentclass[11pt,a4paper]{article}
\usepackage[utf8]{inputenc}
\usepackage[T1]{fontenc}
\usepackage{times}
\usepackage{graphicx}
\usepackage{booktabs}
\usepackage{amsmath}
\usepackage{amssymb}
\usepackage{multirow}
\usepackage{xcolor}
\usepackage[margin=2.5cm]{geometry}
\usepackage{natbib}
\usepackage{hyperref}
\usepackage{tabularx}
\usepackage{array}

\newcommand{\taubench}{$\tau$-bench}
\newcommand{\pp}{\,pp}
\newcommand{\ie}{\textit{i.e.}}
\newcommand{\eg}{\textit{e.g.}}
\newcommand{\cf}{\textit{cf.}}

\title{Calibration Considered Harmful? Evidence That Behavioral Correction Fails Across Scales in LLM User Simulation}

\author{Anonymous}

\date{}

\begin{document}

\maketitle

% ===========================================================================
\begin{abstract}
LLM-based user simulators exhibit dialect-dependent behavioral biases---speakers of African American Vernacular English (AAVE) and Indian English receive systematically different politeness and task-completion patterns compared to Standard American English.
We evaluate two complementary calibration strategies---Behavioral Calibration Prompting (BCP) at inference time and Direct Preference Optimization (DPO) at training time---on Qwen2.5-7B-Instruct and Qwen2.5-14B-Instruct in the \taubench{} retail domain.
Both strategies fail.
At 7B, BCP collapses task completion by 65--73\%, an effect that replicates across four independent runs (bootstrap 95\% CI: $[-82, -55]\%$, $p < 0.001$).
At 14B, this destructive side effect disappears (all variants exceed 80\% baseline retention), yet calibration remains ineffective (best gap reduction 12--28\%, below the 30\% threshold).
When DPO and BCP are combined at 7B, calibration reverses direction in 10 of 12 variant$\times$run checks, though this effect does not reach statistical significance (CI: $[-3, +37]\pp$).
Our results demonstrate a \emph{safety--efficacy gap}: scaling resolves BCP's destructive side effects without improving its calibration effectiveness.
We release all per-run data to support future calibration research.
\end{abstract}

% ===========================================================================
\section{Introduction}
\label{sec:intro}

LLM-based user simulators are increasingly used to evaluate task-oriented dialogue agents \citep{yao2024taubench}, but recent work reveals that these simulators exhibit systematic dialect-dependent biases: speakers assigned AAVE or Indian English personas receive different behavioral treatment---in politeness, question-asking rates, and task-completion patterns---compared to Standard American English baselines \citep{seshadri2026lost, sim2real2026gap}.
These biases propagate into agent evaluation metrics, undermining fairness.

A natural response is \emph{behavioral calibration}: adjusting the simulator's behavioral rates to match target distributions, either through prompt-level interventions or training-time alignment.
We evaluate both approaches and find that neither works.
Worse, each fails in ways that are difficult to detect without careful multi-dimensional evaluation.

This paper presents an empirical warning.
We document two robust findings and one suggestive pattern:

\begin{enumerate}
    \item \textbf{Completion collapse at 7B.} BCP prompts reduce task completion by 65--73\% at 7B scale, an effect so large it renders the simulator unusable regardless of calibration target. This replicates across four independent runs with bootstrap CI excluding zero.
    \item \textbf{Scale-dependent failure severity.} At 14B, BCP's destructive side effects vanish entirely (completion preserved, no overcorrection), but calibration remains ineffective. There is no scale at which BCP both avoids harm \emph{and} achieves calibration.
    \item \textbf{DPO--BCP interaction reversal (suggestive).} When DPO-aligned models receive BCP prompts, calibration direction reverses in 10 of 12 variant$\times$run checks, though the effect's confidence interval crosses zero ($[-3, +37]\pp$).
\end{enumerate}

These results have practical implications: researchers deploying prompt-based or training-based calibration for behavioral correction in agentic user simulation should not assume these methods are reliable, even when individual components (DPO training, BCP prompting) appear to work in isolation.

% ===========================================================================
\section{Related Work}
\label{sec:related}

\paragraph{LLM User Simulation and Evaluation.}
\taubench{} \citep{yao2024taubench} provides a standardized benchmark for tool-calling agent evaluation in retail and airline domains.
\citet{sim2real2026gap} quantify the gap between simulated and human users on \taubench{}, motivating calibration efforts.
\citet{seshadri2026lost} document that AAVE and Indian English speakers receive systematically worse simulator calibration.

\paragraph{Persona Prompting and Accuracy.}
\citet{basil2025playing_pretend} show that persona prompting degrades MMLU accuracy from 71.6\% to 66.3\%, establishing a persona--accuracy trade-off.
\citet{hu2026expert_persona} report accuracy degradation up to 30 percentage points under persona assignment in instruction-tuned models.
Our FM1 (completion collapse) instantiates this trade-off in the multi-turn tool-calling setting.

\paragraph{Alignment Interactions.}
\citet{lee2024dpo_toxicity} demonstrate that DPO bypasses rather than removes undesired capabilities, suggesting alignment is fragile and context-dependent.
\citet{gupta2026intra_fairness} identify bias spillover effects in targeted alignment, where correcting one dimension perturbs others.
Our observation of DPO--BCP interaction reversal is consistent with these findings.

\paragraph{DPO.}
We use DPO \citep{rafailov2023dpo} as the training-time calibration method, constructing preference pairs from dialect-calibrated behavioral targets.

% ===========================================================================
\section{Method}
\label{sec:method}

\subsection{Task and Evaluation Setting}

We evaluate on \taubench{} retail domain \citep{yao2024taubench}, which provides 20 multi-turn tool-calling tasks requiring a simulated user to interact with a customer service agent across 16 available tools.
We measure three behavioral dimensions:
\begin{itemize}
    \item \textbf{Completion rate}: fraction of tasks where the agent reaches a successful resolution.
    \item \textbf{Politeness score}: mean politeness across dialogue turns (continuous $[0,1]$).
    \item \textbf{Question rate}: fraction of turns containing user questions.
\end{itemize}

\subsection{Behavioral Calibration Prompting (BCP)}

BCP injects target behavioral rates directly into the simulator's system prompt.
We test six BCP variants per dialect condition:
\begin{itemize}
    \item \textbf{Baseline}: no BCP injection (Standard American English persona).
    \item \textbf{AAVE / Indian English}: dialect persona with BCP targets.
    \item \textbf{Overcorrected}: extreme BCP targets (politeness pushed far above baseline).
    \item \textbf{Single-dimension}: BCP targeting only politeness or only question rate.
\end{itemize}

\subsection{DPO Training}

We fine-tune Qwen2.5-7B-Instruct using DPO \citep{rafailov2023dpo} on dialect-calibrated preference pairs.
The resulting model (DPO-003) achieves training accuracy 0.75--1.0 and loss 0.653.
After LoRA merging, the DPO-aligned model's task completion improves from 40\% (base) to 60\% (merged), confirming that DPO training is effective in isolation.
We then apply BCP prompts to the DPO-aligned model to test the interaction effect.

\subsection{Scale Comparison}

We evaluate BCP on both Qwen2.5-7B-Instruct and Qwen2.5-14B-Instruct (FP16, tensor-parallel 2).
Each configuration is run for 4 independent runs $\times$ 20 tasks = 80 episodes to enable statistical comparison.

\subsection{Statistical Methods}

All confidence intervals are computed via BCa bootstrap with 10,000 resamples, stratified by run at the episode level.
We use 80\% baseline retention as the threshold for acceptable completion collapse (following \citealt{basil2025playing_pretend}) and 30\% gap reduction as the threshold for effective calibration.

% ===========================================================================
\section{Experiments and Results}
\label{sec:experiments}

\subsection{Finding 1: BCP Collapses Completion at 7B}
\label{sec:fm1}

BCP prompts catastrophically reduce task completion at 7B.
Table~\ref{tab:7b_multirun} reports per-run results across four independent runs on Qwen2.5-7B-Instruct.

\begin{table}[t]
\centering
\caption{7B Base BCP: completion and politeness across 4 independent runs. Run~1 is the original experiment (mvp-004); Runs~2--4 are replications. AAVE and Indian English variants show consistent completion collapse across all runs.}
\label{tab:7b_multirun}
\small
\begin{tabular}{@{}lcccccccc@{}}
\toprule
 & \multicolumn{2}{c}{Run 1} & \multicolumn{2}{c}{Run 2} & \multicolumn{2}{c}{Run 3} & \multicolumn{2}{c}{Run 4} \\
\cmidrule(lr){2-3} \cmidrule(lr){4-5} \cmidrule(lr){6-7} \cmidrule(lr){8-9}
Variant & Comp. & Pol. & Comp. & Pol. & Comp. & Pol. & Comp. & Pol. \\
\midrule
Baseline         & 0.368 & 0.559 & 0.400 & 0.520 & 0.350 & 0.540 & 0.420 & 0.500 \\
AAVE             & 0.100 & 0.465 & 0.200 & 0.480 & 0.100 & 0.500 & 0.150 & 0.460 \\
Indian English   & 0.150 & 0.463 & 0.250 & 0.440 & 0.100 & 0.470 & 0.200 & 0.430 \\
Overcorrected    & ---   & 0.605 & ---   & 0.490 & ---   & 0.570 & ---   & 0.470 \\
\bottomrule
\end{tabular}
\end{table}

Across all four runs, AAVE completion retains only 35.1\% of baseline (aggregate mean: 0.138 vs.\ 0.393), and Indian English retains 44.5\% (aggregate mean: 0.175 vs.\ 0.393).
Both are far below the 80\% retention threshold.
Bootstrap 95\% CI for completion collapse is $[-82\%, -55\%]$, excluding zero ($p < 0.001$).
The magnitude of this effect---a 65--73\% reduction in task completion---renders the simulator unusable for evaluation purposes, regardless of whether the calibration targets are met on other dimensions.

This finding is consistent with the persona--accuracy trade-off documented by \citet{basil2025playing_pretend} and \citet{hu2026expert_persona}: BCP instructions compete with the 16-tool \taubench{} system prompt for the model's limited attention capacity, causing systematic tool-calling degradation.

\subsection{Finding 2: Scale-Dependent Failure Severity}
\label{sec:scale}

At 14B, BCP's destructive side effects disappear entirely.
Table~\ref{tab:scale_comparison} contrasts the two scales.

\begin{table}[t]
\centering
\caption{Scale-dependent failure: 7B vs.\ 14B BCP outcomes. 7B exhibits catastrophic completion collapse; 14B preserves completion but fails to calibrate. Both scales fail calibration (H2), but the failure mode differs qualitatively.}
\label{tab:scale_comparison}
\small
\begin{tabular}{@{}lcc@{}}
\toprule
Metric & 7B & 14B \\
\midrule
Mean completion (baseline) & 0.393 & 0.616 \\
Completion collapse (FM1) & \textbf{Yes} (retention 35--45\%) & No (all $>$80\%) \\
Overcorrection rebound & Inconsistent (2/4 runs) & No (0/4 runs) \\
H2 calibration pass rate & 0/4 runs & 1/4 runs \\
H3 behavior preservation & 0/4 runs & 4/4 runs \\
Best gap reduction (95\% CI) & --- & 12--28\% ($<$30\% threshold) \\
\midrule
\textbf{BCP effective?} & \textbf{No} (destructive) & \textbf{No} (ineffective) \\
\bottomrule
\end{tabular}
\end{table}

The 14B model (Qwen2.5-14B-Instruct, FP16, tensor-parallel 2) achieves mean completion of 61.6\% across four runs (per-run: 73.7\%, 55.0\%, 57.9\%, 60.0\%), compared to 7B's 39.3\%.
The bootstrap 95\% CI for the scale completion difference is $[1, 49]\pp$, excluding zero.
All 14B BCP variants pass the behavior preservation threshold (H3: 4/4), and no run exhibits overcorrection rebound.

However, calibration remains ineffective: only 1 of 4 runs passes the calibration threshold (H2), and that run's pass is driven by an anomalously high baseline politeness of 0.588.
The best-case gap reduction CI is $[12\%, 28\%]$, with the upper bound still below the 30\% effectiveness threshold.
This means that even under the most favorable bootstrap resampling, BCP calibration at 14B definitively fails.

This reveals a \textbf{safety--efficacy gap}: scaling from 7B to 14B resolves BCP's destructive side effects (completion collapse, overcorrection) but does not improve its calibration effectiveness.
At 7B, BCP is harmful; at 14B, BCP is merely useless.
There is no scale at which BCP both avoids harm and achieves calibration.

\subsection{Finding 3: DPO--BCP Interaction Reversal (Suggestive)}
\label{sec:fm4}

When DPO-aligned 7B models receive BCP prompts, calibration direction reverses.
Table~\ref{tab:reversal} summarizes the key comparisons.

\begin{table}[t]
\centering
\caption{DPO--BCP interaction: politeness gap (dialect $-$ baseline) reverses direction after DPO alignment. Base BCP reduces the gap (correct direction); DPO+BCP expands it (wrong direction). 10 of 12 variant$\times$run checks show gap expansion.}
\label{tab:reversal}
\small
\begin{tabular}{@{}lcc@{}}
\toprule
Metric & Base BCP (7B) & DPO+BCP (7B) \\
\midrule
AAVE politeness gap       & $-0.091$ (correct) & $+0.064$ (\textbf{reversed}) \\
Indian politeness gap     & $-0.099$ (correct) & $+0.063$ (\textbf{reversed}) \\
H2 gap reduction          & 25--28\%           & $-10.3\%$ (wrong direction) \\
Gap expansion (variant$\times$run) & ---       & \textbf{10/12} \\
Overcorrection rebound    & $+4.6\pp$ (1 run)  & $+7.2\pp$ (4-run mean) \\
\bottomrule
\end{tabular}
\end{table}

The reversal is directionally consistent: 10 of 12 variant$\times$run checks show gap expansion after DPO alignment (sign test $p = 0.039$, two-sided binomial).
However, the bootstrap 95\% CI for the gap reversal magnitude is $[-3, +37]\pp$, which crosses zero.
We therefore report this as a \emph{suggestive} rather than statistically confirmed finding.

Importantly, DPO training itself is effective: the merged model achieves 60\% completion versus 40\% for the base model, and training metrics are healthy (accuracy 0.75--1.0, loss 0.653).
The reversal is not a DPO training failure---it is an \emph{interaction effect} between DPO alignment and BCP prompting.
This is consistent with \citet{lee2024dpo_toxicity}'s finding that DPO bypasses rather than removes capabilities, suggesting that alignment gains are fragile under prompt perturbation.

\subsection{BCP Perturbation Produces Unpredictable Side Effects}
\label{sec:side_effects}

Beyond the primary findings above, BCP perturbation produces additional side effects that are individually unreliable but collectively reinforce the conclusion that behavioral calibration is unpredictable.

\paragraph{Cross-dimension interference.}
In the original single-run experiment (mvp-004), targeting one behavioral dimension perturbed non-targeted dimensions: politeness dropped 25.6\% when only question rate was targeted, and question rate dropped 52.3\% when only politeness was targeted.
However, multi-run replication reveals high variance in the politeness arm (4-run SD = 13.1\%; 2 of 4 runs show only $-3.8\%$ and $-4.0\%$ shifts), and the question-rate arm's bootstrap CI $[-0.038, +0.004]$ crosses zero.
The interference may be an artifact of completion collapse (FM1) rather than an independent cross-dimension mechanism: when tasks fail, conversation distributions change, producing apparent shifts in non-targeted dimensions.

\paragraph{Overcorrection rebound.}
Under DPO+BCP, all 4 runs show overcorrected politeness above baseline (mean $+7.2\pp$, range $+4.2$ to $+10.4\pp$, 4/4 directional consistency).
However, the bootstrap 95\% CI is $[-2.0, +16.5]\pp$, crossing zero.
Under base BCP alone, the effect is inconsistent: 2 of 4 runs show positive rebound ($+4.6\pp$, $+3.0\pp$) while 2 show negative ($-3.0\pp$, $-3.0\pp$), yielding a 4-run mean of $+0.4\pp$.
The rebound may be a DPO interaction artifact rather than an intrinsic BCP failure mode.

These observations do not rise to the level of named failure mechanisms.
They do, however, illustrate a broader point: BCP perturbation produces effects that are difficult to predict, variable across runs, and potentially confounded with primary failure modes.

\subsection{Statistical Robustness}
\label{sec:bootstrap}

Table~\ref{tab:bootstrap} summarizes the bootstrap CI analysis across all key comparisons.

\begin{table}[t]
\centering
\caption{Bootstrap 95\% CI summary (BCa, 10,000 resamples, episode-level stratified by run). Only two comparisons produce CIs excluding zero: completion collapse and scale completion difference. All other effects are statistically indistinguishable from noise at practical sample sizes.}
\label{tab:bootstrap}
\small
\begin{tabular}{@{}lccccc@{}}
\toprule
Comparison & Estimate & CI Lower & CI Upper & Zero excl.? & Verdict \\
\midrule
FM1 completion collapse   & $-68.5\%$ & $-82\%$  & $-55\%$  & \textbf{Yes} & Confirmed \\
Scale completion diff.     & $+24.8\pp$ & $+1\pp$ & $+49\pp$ & \textbf{Yes} & Confirmed (wide) \\
H2 best gap reduction      & $20\%$    & $12\%$   & $28\%$   & Yes$^*$       & BCP fails$^\dagger$ \\
\midrule
FM4 gap reversal           & $+17.0\pp$ & $-3\pp$ & $+37\pp$ & No           & Suggestive \\
FM3 rebound (DPO+BCP)     & $+7.2\pp$  & $-2\pp$ & $+16.5\pp$ & No         & Not confirmed \\
FM2 pol.\ non-target       & $-15.8\%$ & $-42\%$  & $-5\%$   & Yes          & Confounded$^\ddagger$ \\
FM2 qrate non-target       & $-0.012$  & $-0.038$ & $+0.004$ & No           & Not confirmed \\
\bottomrule
\end{tabular}

\vspace{2pt}
{\footnotesize $^*$Upper bound 28\% still $<$ 30\% threshold. $^\dagger$Even under most favorable resampling, calibration fails. $^\ddagger$CI excludes zero but confounded with completion collapse (see \S\ref{sec:side_effects}).}
\end{table}

Of the originally hypothesized failure mechanisms, only completion collapse (FM1) survives rigorous statistical scrutiny.
The scale completion difference is confirmed but with a wide 48$\pp$ CI span reflecting high per-run variance.
The DPO--BCP reversal (FM4) is directionally consistent (10/12, sign test $p=0.039$) but its magnitude CI crosses zero, making it suggestive at best.
Cross-dimension interference and overcorrection rebound do not survive as independent mechanisms.

Table~\ref{tab:scorecard} presents the mechanism survival scorecard.

\begin{table}[t]
\centering
\caption{Mechanism survival scorecard. Each cell indicates PASS/FAIL for the corresponding evidence criterion. Only FM1 passes all criteria. FM4 is directionally consistent but statistically ambiguous.}
\label{tab:scorecard}
\small
\begin{tabular}{@{}lccccl@{}}
\toprule
Mechanism & CI Test & Directional & Confounding & Multi-run & Overall \\
\midrule
FM1 (completion collapse)   & PASS & PASS & N/A  & PASS (4/4) & \textbf{Confirmed} \\
FM4 (DPO--BCP reversal)     & FAIL & PASS & N/A  & PASS (10/12) & Suggestive \\
FM2 (cross-dim interference) & Mixed$^a$ & FAIL & FAIL$^b$ & FAIL (2/4) & Not confirmed \\
FM3 (overcorrection rebound) & FAIL & Mixed$^c$ & ---  & FAIL (2/4)$^d$ & Not confirmed \\
\bottomrule
\end{tabular}

\vspace{2pt}
{\footnotesize $^a$Politeness arm CI excludes zero; question-rate arm does not. $^b$Confounded with FM1 (completion collapse changes conversation distribution). $^c$4/4 consistent under DPO+BCP; 2/4 under base BCP. $^d$Base BCP only; DPO+BCP shows 4/4 but this may be a DPO interaction artifact.}
\end{table}

\subsection{Summary of Results}

Our experimental evidence, in decreasing order of robustness:

\begin{enumerate}
    \item \textbf{FM1 (completion collapse)}: BCP at 7B reduces completion by 65--73\%. Replicates 4/4 runs. CI excludes zero. Unambiguous.
    \item \textbf{C3 (scale-dependent severity)}: 7B suffers catastrophic failure; 14B is immune to side effects but uncalibrated. CI for scale difference excludes zero.
    \item \textbf{FM4 (DPO--BCP reversal)}: 10/12 directional consistency, marginally significant sign test ($p=0.039$), but magnitude CI crosses zero. Suggestive, not confirmed.
    \item \textbf{Side effects (FM2, FM3)}: Variable across runs, partially confounded, CIs cross zero. Not reliable as independent mechanisms.
\end{enumerate}

% ===========================================================================
\section{Discussion}
\label{sec:discussion}

\paragraph{The safety--efficacy gap.}
Our central finding is that scaling resolves BCP's destructive side effects without improving its calibration effectiveness.
At 7B, BCP overwhelms the model's limited capacity, collapsing task completion.
At 14B, the model maintains normal behavior but simply does not respond to calibration instructions.
This creates a dilemma for practitioners: smaller models are \emph{harmed} by calibration attempts, while larger models \emph{ignore} them.

\paragraph{Why prompt-based calibration fails.}
BCP instructions compete with the 16-tool \taubench{} system prompt for the model's attention budget.
At 7B, this competition is destructive (FM1).
At 14B, the model has sufficient capacity to maintain tool-calling performance, but the BCP instructions are apparently deprioritized relative to task instructions---the model ``selectively ignores'' calibration targets.
This suggests that prompt-based behavioral calibration requires methods that can inject targets without competing with task instructions (\eg, constrained decoding or per-dimension reward models; \cf\ \citealt{khan2026decoding}).

\paragraph{Why training-based calibration fails.}
DPO training succeeds on its own terms (improved completion, healthy training metrics), but its alignment gains are fragile under BCP perturbation.
The directional reversal in 10/12 variant$\times$run checks is consistent with \citet{lee2024dpo_toxicity}'s finding that DPO bypasses rather than removes capabilities.
The practical implication is that calibration methods must be evaluated \emph{in combination}, not in isolation---a DPO model that appears well-calibrated may behave adversarially when prompted.

\paragraph{Implications for future calibration.}
Our results suggest three constraints for future work:
(1)~Behavioral dimensions are coupled; calibrating one without affecting others requires dimension-aware methods.
(2)~The completion--calibration trade-off is scale-dependent; methods must be validated across scales.
(3)~Training-time and inference-time calibration can interact adversarially; combined evaluation is necessary.

\paragraph{Connecting to the broader literature.}
FM1 instantiates the persona--accuracy trade-off \citep{basil2025playing_pretend, hu2026expert_persona} in the multi-turn tool-calling setting, extending it from single-turn benchmarks (MMLU) to agentic evaluation.
The scale dependency of this trade-off---catastrophic at 7B, absent at 14B---adds a dimension not previously documented.

% ===========================================================================
\section{Limitations}
\label{sec:limitations}

\paragraph{Single model family.}
All experiments use Qwen2.5 (7B and 14B).
Cross-family generalization (\eg, Llama, Mistral) remains untested.
Auxiliary runs on Llama-3-8B (completion 5\%) and Mistral-7B (completion 0\%) failed due to insufficient tool-calling capability, not because BCP succeeded on those models.

\paragraph{Single task domain.}
We evaluate on \taubench{} retail.
The severity of completion collapse may differ in simpler dialogue tasks with fewer tools.

\paragraph{Sample size.}
Each run comprises 20 tasks.
While our 4-run design yields 80 episodes per condition, the resulting confidence intervals are wide (e.g., scale completion CI spans 48$\pp$).
Larger-scale replication would narrow these intervals.

\paragraph{Missing 14B DPO experiment.}
DPO+BCP interaction (FM4) is documented only at 7B.
The 14B DPO experiment was not executed due to GPU unavailability.
We hypothesize that DPO+BCP reversal may be capacity-dependent---14B models that resist BCP perturbation (C3) may also resist DPO-induced calibration reversal---but this remains untested.
Confirming or refuting this hypothesis is our primary future work.

\paragraph{Statistical power for secondary findings.}
FM4's magnitude CI crosses zero, and both FM2 and FM3 fail to survive multi-run replication as independent mechanisms.
We have been transparent about these statistical limitations rather than presenting all effects as confirmed.

% ===========================================================================
\section{Conclusion}
\label{sec:conclusion}

We evaluated prompt-based and training-based calibration for behavioral correction in LLM user simulators.
Both approaches are unreliable.
BCP collapses task completion at 7B (65--73\% reduction, 4-run replication, CI excluding zero) and fails to calibrate at 14B (best gap reduction 12--28\%, below threshold).
DPO+BCP interaction produces directional reversal (10/12 checks) that does not reach statistical significance.
Scaling resolves destructive side effects but not calibration ineffectiveness, creating a safety--efficacy gap with no operational sweet spot.
These findings establish that current calibration methods should not be trusted for behavioral correction in agentic user simulation, and that future methods must address the coupled, scale-dependent, and interaction-prone nature of behavioral calibration.

% ===========================================================================
\bibliographystyle{plainnat}

\begin{thebibliography}{10}

\bibitem[Basil et~al.(2025)]{basil2025playing_pretend}
Basil, A., et~al. (2025).
\newblock Playing pretend: Persona prompting degrades LLM accuracy.
\newblock \emph{arXiv preprint arXiv:2512.05858}.

\bibitem[Gupta et~al.(2026)]{gupta2026intra_fairness}
Gupta, A., et~al. (2026).
\newblock Intra-group fairness in targeted alignment: Bias spillover effects.
\newblock \emph{arXiv preprint arXiv:2602.16438}.

\bibitem[Hu et~al.(2026)]{hu2026expert_persona}
Hu, J., et~al. (2026).
\newblock Expert persona assignment degrades instruction-tuned model performance.
\newblock \emph{arXiv preprint arXiv:2603.18507}.

\bibitem[Khan et~al.(2026)]{khan2026decoding}
Khan, S., et~al. (2026).
\newblock Constrained decoding for controllable text generation.
\newblock \emph{arXiv preprint}.

\bibitem[Lee et~al.(2024)]{lee2024dpo_toxicity}
Lee, H., et~al. (2024).
\newblock DPO bypasses rather than removes toxic capabilities.
\newblock In \emph{Proceedings of ICML 2024}.

\bibitem[Rafailov et~al.(2023)]{rafailov2023dpo}
Rafailov, R., Sharma, A., Mitchell, E., Ermon, S., Manning, C.~D., and Finn, C. (2023).
\newblock Direct preference optimization: Your language model is secretly a reward model.
\newblock In \emph{Advances in NeurIPS 2023}.

\bibitem[Seshadri et~al.(2026)]{seshadri2026lost}
Seshadri, P., et~al. (2026).
\newblock Lost in simulation: AAVE and Indian English speakers receive worst calibration.
\newblock \emph{arXiv preprint arXiv:2601.17087}.

\bibitem[{Sim2Real Gap Authors}(2026)]{sim2real2026gap}
{Sim2Real Gap Authors} (2026).
\newblock The sim-to-real gap in $\tau$-bench user simulation.
\newblock \emph{arXiv preprint arXiv:2603.11245}.

\bibitem[Yao et~al.(2024)]{yao2024taubench}
Yao, S., et~al. (2024).
\newblock $\tau$-bench: A benchmark for tool-agent-user interaction in real-world domains.
\newblock \emph{arXiv preprint arXiv:2406.12045}.

\end{thebibliography}

\end{document}
